\section{Strict continuous time}\label{sec:continuous-time}

The strict continuous-time cone \(\Dct\) is a nonempty Euclidean-open subset
of \(\Dplus\).  To see the latter inclusion directly, put

\[
 u=\sqrt{gt/c},\qquad v=\sqrt{ct/g},\qquad w=\sqrt{cg/t}.
\]

Then \(u,v,w\in(0,1)\), \((c,g,t)=(vw,uw,uv)\), and

\[
 1+c-g-t=1+vw-u(v+w)>(1-v)(1-w)>0,
\]

with the other two inequalities cyclic.  Every generic maximal Jacobian minor
has the same rank on \(\Dct\) as on \(\Dplus\), because a nonzero polynomial
cannot vanish on the nonempty open continuous-time parameter domain.

\begin{proof}[Proof of \cref{cor:ct-main}]
Because \(\Dct\) is Euclidean-open in \(\Dplus\), a source-open strict
continuous-time containment witness is also a principal-domain containment
witness.  Necessity therefore follows from \cref{thm:main}; the discussion
below verifies that the constructive and generic ingredients also remain
inside the strict continuous-time domain.

Root movement, true-cut rank, the 204 pointwise obstructions, and the
tree--sunlet sum of squares hold throughout \(\Dplus\), hence throughout
\(\Dct\).  The isotropic slice used for generic noncut recovery contains
strict continuous-time points.  Triple-product marginal sections and
physical local products are supplied by the continuous-time part of
\cref{lem:subdivision,lem:marginal}.  Each polynomial or rank witness remains
nonzero on a dense open subset of \(\Dct\).

The common triangle preimage in \cref{thm:H14} is strict continuous-time.
The simultaneous bridge construction in \cref{lem:gluing} has the uniform
positive actual-edge lower margin \(7\varepsilon/(8U)\), while the effective
isotropic bridge has margin at least \(3\varepsilon/4\).  Thus necessity and sufficiency both
restrict to \(\Dct\), and the same semialgebraic genericity and exact
reconstruction arguments apply.
\end{proof}

No boundary case \(c=gt\), \(g=ct\), \(t=cg\), a zero eigenvalue, or an
inheritance probability in \(\{0,1\}\) is included.
